\begin{resumo}

    A bolsa de valores reflete as mudanças nos preços dos ativos e funciona como um indicador-chave da saúde econômica, evidenciando tendências de alta ou baixa do mercado. Este trabalho mostrou que a previsão desse comportamento evoluiu de modelos clássicos de \textit{Machine Learning} (como SVR, CART e MLP) para arquiteturas de \textit{Deep Learning}, em especial redes LSTM, e modelos de \textit{Ensemble} e \textit{Gradient Boosting}, aplicados tanto à Regressão (previsão de valores) quanto à Classificação (movimento e regimes de risco) em mercados desenvolvidos e emergentes. Além disso, também destacou a relevância da escolha do \textit{dataset}, da granularidade temporal, da engenharia de \textit{features} e, mais recentemente, da Inteligência Artificial Explicável (xAI), sobretudo com o uso do SHAP para explicar as decisões de modelos complexos. Alinhado a esse cenário, este trabalho foca o mercado brasileiro (IBOVESPA), utilizando dados diários de mais de 25 anos obtidos no Yahoo Finance, pré-processados com normalização MinMax e estruturados por meio de \textit{sliding window} de 100 dias, para treinar uma rede LSTM empilhada implementada em Keras/TensorFlow. O desempenho é avaliado por métricas de Regressão (RMSE e $R^2$), enquanto o SHAP é aplicado para explicar a contribuição temporal e das variáveis de entrada nas previsões. O objetivo deste projeto é desenvolver e analisar um modelo de previsão do preço das ações do IBOVESPA baseado em LSTM, complementado por xAI, de forma a produzir previsões acuradas e explicáveis que baseiem decisões de risco em mercados emergentes e voláteis como o mercado acionário brasileiro.

 \vspace{\onelineskip}
    
 \noindent
 \textbf{Palavras-chave}: LSTM, rede neural, previsão de séries temporais, xAI, SHAP, IBOVESPA.
\end{resumo}
